% page 187 of the facsimile (image p-186 of the scan)
% block 167.6 x 242.0 mm ; left margin 34.4 mm
\pgc{12.700mm}{0.000mm}{
\l{\cel{58}179}
\l{}
\l{}
\l{}
\l{\sou{fregato}. (\sou{nav.)}\cel{3}La maxim granda de la milito-navi kun un fer-}
\l{\cel{3}deko o tota baterio, qua, importale, situesas dop la kombato-}
\l{\cel{3}navo. - DEFIRS.}
\l{}
\l{\sou{fremisar}.\cel{2}(\sou{netrans., de, pro)}\cel{2}I. Agitar su bruisetoze. - II.}
\l{\cel{3}Agitar su konvulse. - FS.}
\l{}
\l{\sou{freno}. Mekanismo uzata por lentigar o haltigar la movo di mashino.}
\l{\cel{3}- FIS.}
\l{}
\l{\sou{frenezio}. I. Deliro furioza. - II. (\sou{metaf.)}\cel{2}Eceso en ula pasiono.}
\l{\cel{3}- DEFIS.}
\l{}
\l{\sou{frenologio}. Doktrino segun qua singla de nia fakultati od ins-}
\l{\cel{3}tinti havas kom situeso ta o ca parto di la cerebro, ed esas}
\l{\cel{3}konocebla per la observo di la protuberanci o di la kavaji ko-}
\l{\cel{3}respondanta che la kranio. - DEFIS.}
\l{}
\l{\sou{frequa}. Qua eventas multafoye, tre mikra-intervale. - DEFIS.}
\l{}
\l{\sou{frequenco}. (\sou{elektrotekniko}). (\sou{pri la elektro-flui alternanta}).}
\l{\cel{3}La nombro de ocili dum un sekundo.}
\l{}
\l{frequentar. (\sou{trans.}) I. Venar (a loko) granda-nombre.- II. Venar}
\l{\cel{3}(a loko) ofte. - III. Venar freque (kun persono) kom akompa-}
\l{\cel{3}nanto. - FIS.}
\l{}
\l{fresha. Qua ne esas velkinta. - DEFIRS.}
\l{}
\l{\sou{fresko}. I. Pikturo sur muro, qua, facite per aquo-farbi sur}
\l{\cel{3}induturo tre recenta, sikeskas e hardeskas kun lu. - II.}
\l{\cel{3}Io quo esas piktita per ta procedo. - DEFIRS.}
\l{}
\l{fretar. - (\sou{trans).}\cel{2}Lokacar navo por transportar vari.-DEFS.}
\l{}
\l{frezar. (\sou{trans.}) (\sou{tekn.}) I. - Mi-bizelizar la orifico di truo en}
\l{\cel{3}qua esas insertota skrubo od irgaltra objekto. - II. Laborar,}
\l{\cel{3}entaliar la ligno o la metali per solido rotaciva qua esas}
\l{\cel{3}garnisita ye denti, per qua - kande la solido rotacas - onu po-}
\l{\cel{3}vas deprenar spani de la surfaco kun qua onu igas lu kontak-}
\l{\cel{3}tar tre preso. - F.}
\l{}
\l{friabla. Qua esas facile reduktebla a peceti. - EFIS.}
\l{}
\l{frianda. (Kozo) qua luras per ulo delikata segun quante koncern-}
\l{\cel{3}esas la palato. F.}
\l{}
\l{fricionar. (\sou{trans., per}). Frotar parto di la homo-korpo por igar}
\l{\cel{3}plu rapida la cirkulo di sango, por kuracar doloro, ec. DEFIS.}
\l{}
\l{frigoro. To quo rezultas de la basigo di temperaturo ; absenteso}
\l{\cel{3}kompleta o parta di kaloro. - EF}
}
